%\documentclass{report}
\documentclass{article}
\usepackage{graphicx} % required for inserting images
\usepackage[pagebackref]{hyperref}
\usepackage{xcolor} % urls
\usepackage{fancyhdr} % header
\usepackage{mathptmx} % times font


\title{Tech Writings}
\author{Yasmeen Dao}
\date{August 2026}


\definecolor{DarkBlue}{rgb}{0.0, 0.0, 0.545}
\hypersetup{
	colorlinks=true,         % enable colored links
	linkcolor=DarkBlue,    % color for internal links
    urlcolor = DarkBlue
}
\pagestyle{fancy}
\fancyhf{} % clear default header/footer
\fancyhead[L]{Tech Writings}
%\fancyhead[R]{\leftmark}   % shows current chapter name
\fancyfoot[C]{\thepage}
%\fancyhead[L]{\leftmark}
\fancyhead[R]{\nouppercase{\rightmark}}
\renewcommand{\refname}{Bibliography}


\begin{document}

\maketitle
\tableofcontents 

\clearpage

\section{Harvest Now, Decrypt Later: The Importance of Cryptographic Migration Now}
%\addcontentsline{toc}{chapter}{Harvest Now, Decrypt Later: The Importance of Cryptographic Migration Now}
\markboth{Harvest Now, Decrypt Later
}{Harvest Now, Decrypt Later: The Importance of Cryptographic Migration Now
}
Sept 25, 2026 \\

In spite of the uncertain quantum computing timelines, this assessment looks to treat the quantum threat as a present risk instead of a future one, due to demonstrated adversary capabilities and the potential scale of economic disaster. 

\subsection{Introduction}

With the advent of quantum computers, problems in drug discovery, optimization, and materials science may see solutions that were once computationally infeasible. Despite the many promising applications, one of the negative consequences of quantum computers is breaking current cryptography. Quantum computers can solve the mathematical problems that build the foundation of most public-key cryptography we use today, including Rivest-Shamir-Adleman (RSA), Diffie-Hellman (DH) key exchange, and Elliptic Curve Diffie-Hellman (ECDH) key exchange. Specifically, Shor's algorithm poses a threat to these public-key cryptography schemes by efficiently finding the prime factors of large numbers. For comparison, it would take a modern supercomputer millions of years to break a sufficiently sized RSA key. However, a quantum computer with about 4,000 qubits could take a few hours. On the day when a cryptographically relevant quantum computer (CRQC) can use Shor's algorithm to break these cryptographic schemes, most of the encrypted data sent across computer networks, national security, health records, and financial transactions will be able to be decrypted by bad actors. 

Unlike Y2K, the transition to a four-digit year format, the date of this transition was known. In the case of quantum computers, we don't know when that day will be. Bad actors, or adversaries, can collect encrypted data today for later decryption by a CRQC, with the prevalence of "Harvest Now, Decrypt Later" (HNDL), also known as “Store Now, Decrypt Later” (SNDL), attacks. HNDL attacks are a threat now because most data has a lifetime of more than ten years, and information such as bank records, health records, and national security secrets need privacy for the next 30-plus years. These attacks could allow foreign nation-states and hackers to see today's data decrypted in the next decade. 

\subsection{Threat Analysis}

In their August 2023 fact sheet, the Cybersecurity and Infrastructure Security Agency (CISA), the National Security Agency (NSA), and the National Institute of Standards and Technology (NIST) state, “Early planning is necessary, as cyber threat actors could be targeting data today that would still require protection in the future (or in other words, has a long secrecy lifetime), using a catch now, break later or harvest now, decrypt later operation” [1]. This reflects the urgency from leading cybersecurity authorities that the threat is active rather than theoretical. Furthermore, this assessment aligns with demonstrated threat behavior. Although not described as an HNDL attack, the 2020 SolarWinds cyberattack demonstrated similar capabilities the model depends on. The threat actor, identified as the Russian Foreign Intelligence Service (SVR), gained undetected access to thousands of organizations’ computer networks, including multiple federal agencies [2]. Access at scale, long-term patience, and targeting of high-value infrastructure are the abilities a bad actor would need in an HNDL attack today. 

\subsection {The Migration to Post-Quantum Cryptography}

To defend against attacks from quantum computers, existing encryption algorithms are no longer secure. Organizations must be prepared with post-quantum encryption algorithms. These algorithms contain mathematical problems that are more difficult for both conventional and quantum computers to solve. Two main tasks these algorithms are designed for are general encryption, protecting data exchanged across a public network, and digital signatures, authenticating identity and ensuring data integrity. Still, the migration to post-quantum cryptography (PQC) alone isn't enough. This will likely not be the last cryptographic upgrade, as algorithms and standards are constantly changing and updating to meet new threats. Thus, implementing cryptographic agility, the ability to swap standards and algorithms, is vital. This avoids re-architecting or “ripping and replacing” systems every time there is a new security upgrade. Y2K was averted because software engineers and the government took proactive measures to patch and fix the code before the deadline arrived. The same proactive migration needs to occur with post-quantum cryptography even without the certainty of a fixed date. 

\subsection{Regulatory Insight}

The urgency of post-quantum cryptography continues to be backed by regulatory insight. In 2024, NIST completed an eight-year search for algorithms to use as the standard for post-quantum cryptography (PQC). The three published standards are FIPS 203 ML-KEM for key agreement, and FIPS 204 ML-DSA and FIPS 205 SLH-DSA for digital signatures [3]. The search also gives guidance towards future cybersecurity efforts, which will have to comply with NIST post-quantum cryptography standards. Additionally, NSA’s CNSA 2.0 suite states that starting January 1, 2027, no new technology that’s bought for national security applications will be allowed that doesn’t support PQC [4]. This creates a hard deadline that not just government agencies, but also contractors and vendors, must meet. On June 22, 2026, Executive Order 14412 was signed, giving agencies until December 31, 2030, to move high-impact systems and assets to post-quantum key exchanges and until December 31, 2031, for signatures [5]. The order reinforces the urgency of the HNDL threat model, stating that “Ongoing cyber activity against our Nation also presents the risk of adversaries collecting United States information now, and decrypting it later once large-scale quantum computers are operational” [5]. Furthermore, this order acknowledges that the migration is not only a domestic effort. The Secretary of State should work with NIST, the Department of Homeland Security (DHS), and intelligence communities to “identify and engage foreign governments and industry groups in key countries to encourage their transition to PQC algorithms standardized by NIST” [5]. This reflects a broader understanding that international coordination of PQC migration matters. Even if a nation completes its own migration, its dependence on international commercial and financial networks leaves it vulnerable if those networks remain unmigrated or outdated. 

\subsection{Uncertainty Argument}

A major uncertainty in assessing this threat is timing. Estimates for when a CRQC will exist vary widely. Companies like Google and Cloudflare have set their PQC migration deadline to 2029, while the EU has set a 2030 target for securing high-risk critical information infrastructure and 2035 for medium risk [6]. Meanwhile, error-corrected qubits at scale present hardware barriers.
IBM is committed to 200 logical qubits by 2029. Other estimates suggest the 2030s to 2040s for a CRQC, but depend on hardware and algorithm progression. Because of these estimations, security leaders, CISOs, or CTOs might question the necessity of putting care, budget, and time into migration now when they currently have immediate and high-priority tasks to deal with. It could even be possible that by the time a CRQC exists, some data such as routine credit card transactions or old operational data may have already lost its sensitivity. However, leaders in security are viewing quantum readiness as a hardware deployment issue, when it is an immediate data exposure and continuous management issue. State intelligence and defense, intellectual property, and health records still have confidentiality requirements that are measured in decades. Long-retention data is a real risk. This analysis treats the quantum threat as a current issue not because the timeline is set in stone, but because the cost of migrating too late can be catastrophic. A 2023 Hudson Institute study, cited in a Citi Bank report, estimates that just one quantum attack on one of the top five major banks could put 2.0 to 3.3 trillion of U.S. GDP at risk [7]. Investing in system protection now can help mitigate a far greater financial cost later. 

\subsection{Recommendations}

An important note to make is that migrations take time. Organizations and security leaders often underestimate the amount of time it takes to transition all assets to post-quantum. Many current systems have legacy infrastructure with outdated and unorganized algorithms and protocols. Therefore, the sooner organizations start migrating their cryptography, the less risk they will face later. 

Migration should begin with clear ownership that is responsible for a prioritized migration plan. From there, leads should first start inventorying their data and security systems to prioritize critical infrastructure and data lifespan. These high-value assets, or systems and critical information that carry the most sensitive data and risk, are critical. They need to know what algorithms are in use, where they are located, and what is exposed to the quantum threat. During the transition, organizations should work to implement a hybrid post-quantum cryptography strategy that utilizes classical encryption methods with quantum-resistant algorithms. This is important to protect against both existing threats and the possibility that a new PQC algorithm has a weakness itself, as it is less tested than older classical algorithms. Lastly, the continuous monitoring of this inventory is equally important to keep up with changes and risks. One-time assessments leave the same gaps that allowed the SolarWinds attack to remain undetected for months [2]. 

\subsection{Looking Forward}

Without encryption that is quantum-resilient, organizations are vulnerable to HNDL attacks now, as well as potential catastrophic data breaches in the future. Government deadlines continue to be set, while several companies have moved their migration targets to earlier, signaling the shift from theoretical concern to practical urgency. Observing the logical qubit and error correction breakthroughs in hardware, as indicators, will help signal that the CRQC timeline is accelerating. This provides a tangible sense of urgency in migration. Nonetheless, organizations that wait for certainty are choosing to wait until the risk is already realized, at which point it might be too late. 



% before the bibliography
\clearpage
\phantomsection
\addcontentsline{toc}{section}{Bibliography}
%\addcontentsline{toc}{chapter}{Bibliography}
\markboth{Bibliography
} {Bibliography}
\begin{thebibliography}{7}

\bibitem{cisa2023}
Cybersecurity and Infrastructure Security Agency, ``CISA, NSA, and NIST Publish Factsheet on Quantum Readiness,'' \textit{CISA News}, Aug. 21, 2023. [Online]. Available: \url{https://www.cisa.gov/news-events/alerts/2023/08/21/cisa-nsa-and-nist-publish-factsheet-quantum-readiness}

\bibitem{cisa2021}
Cybersecurity and Infrastructure Security Agency, ``Joint NCSC-CISA-FBI-NSA Cybersecurity Advisory on Russian SVR Activity,'' \textit{CISA Alert}, May 7, 2021. [Online]. Available: \url{https://www.cisa.gov/news-events/alerts/2021/05/07/joint-ncsc-cisa-fbi-nsa-cybersecurity-advisory-russian-svr-activity}

\bibitem{nist2024}
National Institute of Standards and Technology, ``NIST Releases First 3 Finalized Post-Quantum Encryption Standards,'' \textit{ NIST News}, Aug. 13, 2024. [Online]. Available: \url{https://www.nist.gov/news-events/news/2024/08/nist-releases-first-3-finalized-post-quantum-encryption-standards}

\bibitem{nsa2022}
National Security Agency, ``Commercial National Security Algorithm Suite 2.0 (CNSA 2.0),'' \textit{ Cybersecurity Advisory, NSA}, Sept. 7, 2022 (updated May 30, 2025). [Online]. Available: \url{https://media.defense.gov/2025/May/30/2003728741/-1/-1/0/CSA_CNSA_2.0_ALGORITHMS.PDF}


\bibitem{eo14412}
Exec. Order No. 14412, ``Securing the Nation Against Advanced Cryptographic Attacks,'' 91 Fed. Reg. 38483, Jun. 25, 2026. [Online]. Available: \url{https://www.federalregister.gov/documents/2026/06/25/2026-12909/securing-the-nation-against-advanced-cryptographic-attacks}

\bibitem{insider2026}
M. Ur Rehman, ``Post-Quantum Cryptography Timelines: When Will Organizations Migrate?,'' \textit{The Quantum Insider}, Aug. 7, 2026. [Online]. Available: \url{https://thequantuminsider.com/2026/08/07/post-quantum-cryptography-timelines/}

\bibitem{citi2026}
S. Bantanidis, K. Master, R. Shah, P. Sharma, and R. Ghose, ``Quantum Threat: The Trillion-Dollar Security Race Is On,'' Citi Institute, Jan. 2026. [Online]. Available: \url{https://www.citigroup.com/rcs/citigpa/storage/public/Citi_Institute_Quantum_Threat.pdf}

\end{thebibliography}

\clearpage
\section{AI, Quantum, Ethics and Relation to Data Security}
%\addcontentsline{toc}{chapter}{AI, Quantum, Ethics and Relation to Data Security}
\markboth{AI, Quantum, Ethics and Relation to Data Security}{AI, Quantum, Ethics and Relation to Data Security}
Draft - 2026
\subsection{Abstract}

Pain point: Today the buzz is about AI. Companies are integrating AI to keep up with technological upgrades. People use it for the convenience of tasks. But the foundation is data, and equally as important is the security of this data. The nexus of AI and quantum is a powerful advancement that can revolutionize industries, but before we get lost in the product growth, we must bring to attention the protection of this data. 

\subsection{Introduction}
Fundamentally, an agent is something that perceives and acts autonomously upon its environment. AI agents are everywhere and in numerous forms. They enhance large language models and generative AI (create images, text, and video from human language interaction). From Alexa to online shopping agents, these computer programs can autonomously make decisions and perform tasks to achieve specific goals. Agentic AI takes it a step further, incorporating multiple agents. They offer advanced autonomy by setting their own goals, planning, and adapting to achieve tasks. The foundation of AI agents (and AI broadly) is data. Most of which is our data. Every day, these programs integrate further into our lives. However, as autonomy increases, so do potential vulnerabilities. Failure to fully understand the capabilities means we are missing out on both the critical governance and security challenges, as well as the true potential of these rapidly evolving technologies.\\

\subsection{Data Security and Privacy Before Product}

AI gets smarter with data. More data means a larger database and increased exposure. As we depend more on AI agents, convenience starts being traded for privacy. To minimize this, it is important to understand the basic functionalities and ensure user needs before advanced transformations can be fulfilled. Maslow’s hierarchy of digital needs is a good blueprint to highlight gaps in current digital foundations. Data security is one of the base-level requirements of Maslow’s hierarchy framework. Before pushing advancements, protecting personal data, maintaining system stability, and ensuring security are critical. Yes, AI agents are helpful, but the data and infrastructure need to be stable and secure first. 

\subsubsection{What Data is at Risk When Sharing With Agents} 

Data at risk can include stored enterprise, personal, federal, and medical data. Additionally, data in transit, like communications.

\subsubsection{The Importance of Protecting Data Today}
What does quantum computing mean for AI and data security? Quantum enhances the processing power of AI, and AI can help quantum systems scale. However, without strong data management, the data powering both remains at risk. Two avenues to explore are: data protection and data analysis. 

\textbf{Data Protection:} The emergence of quantum computers leads to security challenges. The data that travels over networks can be vulnerable today due to SNDL attacks. Data can be harvested now for decryption by a cryptographically relevant quantum computer later. If data is not encrypted with quantum-safe encryption, bad actors can gain access. Therefore, adopting post-quantum cryptography is imperative. Furthermore, incorporating a crypto-agile solution is essential to adapt to changing algorithms and standards without rebuilding system infrastructure for every new update.   

\textbf{Data Analysis:} Another avenue is the analysis of data. Quantum computers can look at the data you share and analyze it in ways we can’t currently do with today’s tools. (complex and different types of problems) However, similar to AI, the analysis of data that you share has the possibility of being analyzed in unintended ways. (Ethical risk) 

\subsection{Responsible Innovation}

AI and agents will continue to be used. They aren’t going away. Similarly, data will continue to be available to a lot of people, whether stored, traveling, or being used by other companies. Especially as quantum computers emerge, post-quantum security is vital for AI to continue safely. Being aware of and understanding the risks is a key part of the solution for securing data in transit and at rest. It is important to make sure federal governments and organizations are responsible for securing data with quantum-safe encryption. Equally important, we, as consumers, must be aware of how our data is being used securely, because it is our data at stake. Data cannot just be a commodity. It is an asset to organizations, but to people, it’s their identity, and in the face of attacks, it can be a liability. 

Data security is the base-level requirement in the digital hierarchy; therefore, post-quantum cryptography is the needed reinforcement. Our digital foundations need to be stronger than the technologies built upon them. 



\clearpage
%/////////////proy
\section{Quantum-Safe Proxy - Network Orchestration, Cryptographic Debt}
%\addcontentsline{toc}{chapter}{Quantum-Safe Proxy - Network Orchestration, Cryptographic Debt}
\markboth{Quantum-Safe Proxy}{Quantum-Safe Proxy}

Draft - 2026
\subsection{Abstract}

Pain point: (Migration) changes or hardware revamping, crypto-agility. Allows for updating of security policies, keys, TLS without code changes to the application. 
Helps with eliminating cryptographic debt. 
Targeting general and security architects

Migration to quantum security may seem overwhelming. How do organizations choose the best strategy? Architects and security leaders to select proper building blocks and define a migration path for organizations to address their system architectures. It is crucial for organizations to choose an agile architectural pattern for their environment to future-proof applications and systems against cryptographic changes. Although doing so may lead to short-term costs, it leads to long-term simplicity in upgrading. 
\subsection{Main}
\textbf{The problem:}
Organizations might look towards replacing an application that uses quantum-vulnerable encryption such as RSA with a post-quantum algorithm. However, this only proves to be an immediate solution. Additional problems:
\begin{itemize}
    \item Revamping applications one at a time is not efficient
    \item Time-consuming (months/years for each application and having to do that for every application
    \item Expensive
    \item Very difficult to map all cryptographic dependencies, as many assets are deep or hidden, leading to incompleteness 
    \item Not future-proof and adaptable. When there are new standards or threats, the process restarts
\end{itemize}

\subsection{What Is a Proxy? (cryptographic agility component)}
One of the foundational building blocks of quantum security (alongside cryptographic libraries) are proxies. A proxy is a router or system that provides a gateway between users and the internet. By filtering content or balancing network loads, proxies act as a middleman or bridge to enable secure communication. To add post-quantum support, proxies layer on existing servers or services. Throughout the environment, orchestration layers help manage the proxies. (man in the middle attacks?) 

\subsubsection{How Does Proxy-Based Network Orchestration Eliminate Cryptographic Debt}
Cryptographic debt is the accumulation of operational inefficiency and risk caused by outdated ciphers, expired certificates, hard-coded cryptography, and fragile key lifecycles that are buried in applications. To eliminate this, it is vital to centralize the cryptographic control at the network layer instead of fixing cryptography in apps one by one. The pros of this: 
\begin{itemize}
    \item Proxy layer to existing service, application, or system
    \item System supports latest cryptographic algorithms through the proxy so no change to code
    \item Upgrade algorithms, TLS versions, policies, keys etc. without rewriting changes to the application
    \item Centralized governance keys, certificates, and policies are all on one control plane
    \item Implementing crypto-agility with the ability to swap algorithms as new threats and standards arrive
\item Implementing post-quantum cryptography without disruption 
\item Ability to scale
\end{itemize}

This approach allows for minimizing rip-and-replace with the goal of achieving cryptographic agility across security tech stacks and environments. This approach also replaces static encryption components with proxy-based options that enable agility. 

\clearpage
%/////// Learnings
\section{Learning}
% \addcontentsline{toc}{chapter}{Learning}
\markboth{Learning 
} {Learning
}
\begin{itemize}
    \item Quantum vs classical modes of transportation ship vs car. 
    \item  Citi estimates that a quantum-enabled cyberattack on U.S. bank's access to Fedwire could put 2.0–3.3 trillion of U.S. GDP at risk.
    \item Future large-scale quantum computers can be more energy efficient because computation inside quantum processors uses zero energy. Quantum logic (quantum gates) is reversible and unitary. Energy is spent on IO systems and cooling. Massive classical supercomputers burn tens of megawatts. 
    \item Q: Explain why quantum entanglement creates correlations that can't simply be explained by giving two particles matching instructions in advance. Two gloves in a box eg: Measuring one particle doesn't send a signal to the other. This is why entanglement doesn't allow faster-than-light communication? 
\end{itemize}


\subsection{Case: Disease }
Simulating molecular behavior, allowing for highly targeted cures - Alzheimer's...cancer. Modeling drug body interactions. 
\subsection{Case: Food Insecurity}
Oct 31, 2026 \\
There is more than enough food produced in the world to feed everyone. How quantum computing can help: 
\begin{itemize}
    \item Route Problem: Quantum optimization algorithms for rerouting in case of climate disasters or blockages. 
    \item Food waste: Quantum forecasting and digital twins (virtual twins of supply chains) usage logistics networks can balance production.
    \item Simulation at the molecular level: The environmental and financial cost of growing food is tied to the Haber-Bosch process (a chemical method that fuses nitrogen and hydrogen to make ammonia for fertilizer). Requires massive heat and pressure, consuming nearly 2 percent of all global energy and producing up to 3 percent of all carbon emissions. 
    
    Nitrogenase is a specialized enzyme in nature responsible for nitrogen fixation (process of converting atmospheric nitrogen gas (N2) into ammonia (NH3). These bacteria use the enzyme cluster (FeMoco - iron, molybdenum, cofactor) in Nitrogenase to pull nitrogen out of the air at room temperature. Interactions with this molecule are complex and difficult to replicate on classical computers. Quantum hardware natively maps molecular bonds. When a quantum simulation maps FeMoco electron behavior, chemical engineers can then mass-produce artificial catalysts that mimic nature. Eliminating the need for fossil-fuel-driven fertilizer production and makes the base cost of growing crops cheaper. 

\end{itemize}

The current bottleneck is in "engineering". Building machines large enough to sustain millions of error-corrected qubits without collapsing.  

\subsection{Case: Quantum Sensors/Quantum Security}
Nov. 11 2026 \\
Research shows diamond-based quantum sensor can pick up electromagnetic signals from chips running lattice-based encryption. Computer chips give off signals internal electrical signals to process data and external electromagnetic signals (radiation and noise) from operation. "whispers" Byproducts are side-channels. Defense measures - masking, shielding, noise, rotating keys. Lattice is repeating grid of points in multi-dimensional space, start at origin and add integer multiples of basic direction vectors. Good basis short and perpendicular; bad basis long and vectors all nearly the same direction. 

\end{document}
